\scanpage{555}
\mspage{466}

un tel exemple, où \deleted{\unsure{la|le}} non la construction du couple $(\unsure{\mathcal M|M},\unsure{\mathcal N|N})$, mais la démonstration du fait que $\mathcal N$ n'est pas accessible dans $\mathcal M$, fait appel à des hypothèses particulières
sur l'univers ambiant, ne doit pas être considéré comme (tout à fait !) ``logiquement accessible''. Pour être vraiment probant, il faudrait un exemple de nature plus élémentaire.

Il reste la possibilité qu'il y ait un th. métamathématique qui assure que \unsure{tt|tout} \unsure{couple} $(\mathcal M,\mathcal N)$ logiquement accessible, avec
% slanted
\margin{C'est faux, \unsure{ex.|en|car} prendre une sous-catégorie fermée, cf \unsure{4.16.7.10|4.16.7 et 10|4.16.5.10|4.16.10} ci-dessous}
$\mathcal M$ accessible, et $\mathcal N$ sous-catégorie str. pl. $L_\pi$-stable dans $\mathcal M$ pour $\pi$ assez grand, soit accessible dans $\mathcal M$. Ceci s'appliquerait en particulier au cas d'une sous-catégorie \uline{\deleted{\unsure{ouverte}}\unsure{fermée|ouverte}} : toute sous-catégorie \deleted{fermée} \deleted{\unsure{de|str.}} (qui soit logiquement accessible) serait accessible dans $\mathcal M$.
% slanted
\margin{\deleted{\unsure{On|Cf} \illegible{2} sous-catégorie fermée, cf \unsure{4.16.7.1|4.16.7} \illegible{1}}}
L'exemple\inserted{-test} le plus simple serait celui du complémentaire \inserted{\unsure{fermé}} $\mathcal M'_z$ de la sous-catégorie ouverte $\mathcal M_z$ ci-dessus$^{*}$, i.e. la catégorie des
% slanted
\margin{NB elle est \unsure{acc.} dans $\mathcal M$ ssi $\mathcal M_z$ l'est !}
$x\in\mathrm{Ob}\,\mathcal M$ tels que $\mathrm{Hom}_{\mathcal M}(x,z)=\varnothing$, \unsure{ou|on} \unsure{non|encore|mor|un}
% archive sheet number: 554
